% !TEX root = ../main.tex
\section{ELF objective and conditioning details}
\label{app:methodobjective}

\paragraph{Time and self-conditioning.}
For a flow row, sample $t=\operatorname{sigmoid}(-1.5+0.8\xi)$ with $\xi\sim\mathcal N(0,1)$ and use noise scale $\sigma=2$. A detached bootstrap pass with zero answer self-conditioning produces $\widehat x^{\,0}$. With $u\sim\mathrm{Bernoulli}(0.5)$, the main pass receives $u\widehat x^{\,0}$ in answer self-conditioning positions. Prompt states are restored to the clean condition independently of $u$. The velocity conversion for both prediction and target uses $\Delta_t=\max(1-t,0.05)$, so the executed target differs from $x-2\epsilon$ where the clamp is active.

\paragraph{SCCFG target.}
Draw the guidance scale $g$ from $\log(1+g)\sim\mathrm{Uniform}(\log1.5,\log6)$, giving $g\in[0.5,5]$. Let $v^0$ and $v^1$ denote auxiliary velocity predictions without and with answer self-conditioning. The main loss uses
\begin{equation}
 \widetilde v^*=\sg\!\left[v^*+u(1-g^{-1})(v^1-v^0)\right].
 \label{eq:methodscguidance}
\end{equation}
The question is present in both auxiliary branches. All auxiliary predictions and targets are detached; during joint training, gradients enter the prompt encoder through the main conditioned forward. At inference, the preceding denoising prediction supplies recurrent self-conditioning and the requested guidance scale is passed to the network. Guidance scale one retains recurrent self-conditioning. Prompt-drop probability is zero in supervised training, and all reported evaluations use prompt CFG one.

\paragraph{Decoder corruption and mixed objective.}
Each row independently draws $b_r\sim\mathrm{Bernoulli}(0.2)$ to choose decoder mode. At each position of a decoder row, sample
\[
 \lambda_{rj}=\operatorname{sigmoid}(0.8+0.8\xi_{rj}),\qquad
 z^{\mathrm{dec}}_{rj}=\lambda_{rj}x_{rj}+(1-\lambda_{rj})\epsilon'_{rj},
\]
where $\xi_{rj}$ and $\epsilon'_{rj}$ are standard Gaussian. Decoder noise scale is one. The model receives decoder mode, time one, zero answer self-conditioning, and the clean question prefix. Token CE predicts the original token at each position, without a next-token shift.

For valid response mask $M_{rj}$, including the terminal token and excluding question/padding positions, a rank-local microbatch minimizes
\begin{equation}
 \mathcal L_{\mathrm{ELF}}^{\mathrm{micro}}=
 \frac{\sum_{rj}M_{rj}\left[
 (1-b_r)\|v_{\theta,rj}-\widetilde v^*_{rj}\|_2^2/d
 -b_r\log p_{\theta,rj}^{\mathrm{dec}}(y_{rj})\right]}
 {\sum_{rj}M_{rj}}.
 \label{eq:methodmicroobjective}
\end{equation}
There is one denominator across both row types, with no additional normalization by mode probability. The flow term averages all 1,024 coordinates, preserving the relative 256/256/512 group widths. Gradients are averaged over data-parallel ranks and accumulation microsteps. ELF-B uses 12 bidirectional blocks of width 768 and ELF-L uses 32 blocks of width 1,280; both share their backbone between flow and token-decoder modes.

\paragraph{Prompt encoder.}
The frozen Qwen token table maps vocabulary IDs to 2,560-dimensional vectors. The learned frontend applies $2560\to512\to d_{\mathrm{model}}$ linear maps without an intervening nonlinearity. Six bidirectional Transformer blocks use width 768 and 12 heads for ELF-B, or width 1,280 and 16 heads for ELF-L, followed by RMSNorm and a 1,024-dimensional output projection. Inputs contain only the question prefix, and padding is masked. The standalone objective averages MSE over positions and coordinates within each question, then weights questions equally. Joint training retains the flow/decoder objective and fixed answer targets. The token table stays frozen while the prompt Transformer and ELF receive gradients from the main forward.

\paragraph{Training-curve cohorts.}
Figure~\ref{fig:promptcurriculum} uses the same two-seed pre-NFT evaluations as the original training-curve campaign: epochs 4, 8, 12 before prompt substitution; epoch 12 immediately after substitution; and epochs 16/18 for ELF-B or 16/20/21 for ELF-L. All 20 plotted points retain their saved predictions, canonical scores, and 4,000-replicate problem-bootstrap intervals. Both epoch-12 conditions share identical ELF weights. These curves use 32 steps and SCCFG 2, whereas the headline results use 64 steps and SCCFG 3 for MATH500. The joint-from-start comparison in Figure~\ref{fig:representationcomparison} instead uses the common $t=\tau^2$ clock and ELF EMA $0.999$ for its six-epoch controls; its four-seed endpoint means are not points on the curriculum curves.

\section{Inference clocks and reconstruction diagnostics}
\label{app:inferenceclocks}

\paragraph{Sampling grid and conditioning.}
For $S$ denoising steps, we set $\tau_0=0$, $\tau_S=1$, and use the same logit-normal-quantile grid across schedules:
\begin{equation}
 \tau_k=\operatorname{sigmoid}\!\left[-1.5+0.8\Phi^{-1}(k/S)\right],
 \qquad 0<k<S.
\end{equation}
We retain the full $f'_\ell(\tau)$ factor in Eq.~\eqref{eq:asyncsampler}. The three time embedders are identical copies of the selected scalar-time checkpoint weights; all other tensors are unchanged. Their outputs are weighted equally, irrespective of the 256/256/512 group widths, with the mean accumulated in FP64 before casting back to model dtype. The resulting embedding is a shared network condition, not a separate Transformer for each group. In particular, permuting the three exponents preserves this conditioning average while changing the coordinate groups to which the local velocity factors apply. Because the group widths differ, that permutation need not preserve total latent noise energy.

The historical single-clock control holds physical powers $(2.5,2,1.5)$, Euler derivatives, 32-step grid, and seed 42 fixed. Replacing $[e(\tau^{2.5})+e(\tau^2)+e(\tau^{1.5})]/3$ by $e((t_{16}t_{24}t_{32})^{1/3})=e(\tau^2)$ reduces ELF-B epoch-18 GSM8K accuracy from 508/1,319 (38.51\%) to 497/1,319 (37.68\%). This is a single-seed comparison against a \emph{geometric}-mean scalar clock; it does not test embedding the arithmetic mean of the times. It uses the historical ELF/prompt EMA $0.9999/0.999$ and native-dtype embedding reduction, so we keep it separate from the later matched tables.

\paragraph{Schedule evidence and operating-point dependence.}
Table~\ref{tab:inferenceclockfull} reports all six schedules in the matched epoch-21 development and test sweeps. The MATH development result corroborates the headline allocation retrospectively; it was not the original selection experiment. The highest GSM8K development mean instead uses $(1.5,2,2.5)$. On the test set, the two additional shared-power schedules achieve higher MATH500 point estimates than the four schedules in the main table. Thus the main comparison supports the usefulness of time reparameterization and the chosen operating point, without identifying a universally optimal exponent vector.

% !TEX root = ../main.tex
\begin{table}[htbp]
\centering\small
\begin{tabular}{lrrrr}
\toprule
& \multicolumn{2}{c}{Development} & \multicolumn{2}{c}{Test} \\
\cmidrule(lr){2-3}\cmidrule(lr){4-5}
Clock powers & GSM8K & MATH & GSM8K & MATH500 \\
\midrule
$(1,1,1)$ & 65.00 & 15.40 & 52.96 & 15.50 \\
$(1.5,1.5,1.5)$ & 71.75 & 17.70 & 57.35 & \textbf{18.10} \\
$(2,2,2)$ & 72.70 & 17.80 & 57.51 & 17.40 \\
$(2.5,2.5,2.5)$ & 71.35 & 16.30 & 55.80 & \textbf{18.10} \\
$(1.5,2,2.5)$ & \textbf{72.80} & 18.30 & 58.07 & 17.60 \\
$(2.5,2,1.5)$ & 72.15 & \textbf{18.70} & \textbf{59.25} & 17.90 \\
\bottomrule
\end{tabular}
\caption{Complete six-schedule ELF-L epoch-21 sweep: mean pass@1 (\%) over seeds 42/123 at 32 denoising steps and SCCFG 2. Development sets contain 1,000 GSM8K and 500 MATH problems. Test columns use the same cohort as Table~\ref{tab:inferenceclocks}. Bold marks each column maximum, including ties.}
\label{tab:inferenceclockfull}
\end{table}


For the 32-step test comparison, async $(2.5,2,1.5)$ minus sync $(2,2,2)$ gives +1.74 percentage points on GSM8K (paired 95\% CI $[+0.72,+2.73]$) and +0.50 on MATH500 ($[-0.90,+1.90]$). Repeating the same two-seed comparison at 16 steps gives 55.57\% versus 54.74\% on GSM8K, a difference of +0.83 points ($[-0.15,+1.82]$), and 16.60\% versus 16.90\% on MATH500, a difference of $-0.30$ points ($[-2.00,+1.40]$). All use SCCFG 2. These are pointwise paired problem-bootstrap intervals with 4,000 replicates, conditional on the fixed checkpoints and generation seeds; the standard deviations in the main table instead describe dispersion across full-benchmark seed accuracies.

\paragraph{Reconstruction intervention.}
Figure~\ref{fig:layerreconstruction} uses ELF-B epoch 18, ELF/prompt EMA $0.9999/0.9999$, SCCFG 1, noise seeds 42/123, and a fixed 64-problem GSM8K development panel. For a selected group $\ell$, we set $t_\ell=2/[2+\rho\,\mathrm{RMS}(x^{(\ell)})]$ and replace its assistant-content coordinates by $t_\ell x^{(\ell)}+2(1-t_\ell)\epsilon^{(\ell)}$, for $\rho\in\{0.8,1.6,2,3\}$. Other coordinates, question positions, terminal tokens, and padding remain clean. The model predicts the full clean latent, which its own token head decodes. The second pass reuses the same corrupted input with the first prediction as self-conditioning; it is not another integration step. CE averages gold-token losses within each problem, then noise seeds and problems equally. Negative changes mean lower CE than the clean-input \emph{prediction} baseline, not a negative loss.

The random control uses eight fixed NumPy PCG64 partitions (seeds 20260918--20260925). Each control slot retains its true-group time and injected Gaussian energy; random-subset RMS is recorded but does not recalibrate the intervention. At $\rho=3$, after self-conditioning, the true-group layer-16 minus layer-24 CE contrast exceeds the corresponding mean-random contrast by 0.967 nats/token (paired 95\% CI $[0.168,1.890]$). The layer-16 minus layer-32 excess is 1.149 ($[0.061,2.241]$), with weaker evidence at lower noise and in the first prediction. Random partitions are averaged within problem before resampling; they are not additional independent problems.

\paragraph{Random-coordinate clocks during generation.}
A separate ELF-L epoch-21 GSM8K control assigns $(2.5,2,1.5)$ to a random 256/256/512 partition for each of seeds 42/123/456/789, keeping that partition fixed throughout generation. At 32 steps and SCCFG 2, actual layer groups yield 58.53\% mean pass@1 and random groups 58.40\%; random minus actual is $-0.13$ points (paired 95\% CI $[-0.80,+0.51]$). This comparison does not demonstrate an accuracy advantage from matching the clocks to the teacher-layer boundaries. Together, the probes establish a distinction between measured local reconstruction geometry and a causal explanation of end-to-end sampling gains.
